\section{Literature Study}
\label{sec:literature}

\subsection{Activity labels and free phrases}

Activity recognizers take RGB, skeleton, IMU, radar, or Wi-Fi, and usually
return one label from a closed set. An open-set model can reject an unknown
action. It still does not describe that action. An open-vocabulary model
can match a clip to a class name while the action is still going on.
Free-form captions also exist. OV-HAR turns a sensor window into a sentence
\cite{ray2025ovhar}, and OpenHOUSE calls a vision-language model for a
free-form description of streaming video \cite{kang2025openhouse}. Both are
evaluated on recorded data, not on a live phone link.

Recent video benchmarks such as Kinetics-400 are clips of about ten
seconds. A single decoded frame, or three frames a second apart, does not
support a claim on that benchmark. Online action detection reads only past
frames and emits a label at every time step. The prototype is causal, in
that it never reads a frame that has not been captured, but it updates only
as fast as the phone model finishes one still image. The missing piece in
the cited caption systems is a sentence produced on the phone, before any
picture is uploaded, while the person is still moving.

\subsection{Retrieval and generative vision-language models}

CLIP retrieves a sentence from a fixed bank by cosine similarity between an
image embedding, or a text embedding, and the bank
\cite{radford2021learning}. It does not write a new sentence. That
retrieval was the checker in the proposal: the phone phrase and the
uploaded picture were each mapped to one bank sentence, and the screen
showed the bank sentence. A generative vision-language model writes the
sentence token by token, and accuracy follows model size. Qwen2-VL scales
this family up to 72B parameters \cite{wang2024qwen2vl}. Models that can
sit on a phone also exist. MiniCPM-V is in the 2B--8B range
\cite{yao2024minicpm}, and SmolVLM-256M fits in under 1~GB
\cite{marafioti2025smolvlm}. Published scores for these models cover
documents, charts, and general questions, not fine differences between
similar activities.

A small generative model can therefore write the sentence on the phone. It
does not, by itself, say which of those sentences are safe to leave on the
screen. The checker in this report follows the text half of the CLIP
comparison and not the image encoder. Both the phone phrase and the laptop
phrase are encoded by a text embedding model, and the score is their cosine
\cite{xiao2024cpack}. Mapping either phrase onto a class sentence for the
still-image evaluation uses the same score:
\begin{equation}
\label{eq:class}
c
=
\arg\max_{c\in\mathcal{C}}
\cos\big(f_T(s), e_c\big).
\end{equation}
The live screen does not apply \eqref{eq:class}. It shows $s$ itself.

\subsection{Client--server collaboration}

Neurosurgeon profiles network layers and chooses where to split them, so
the server receives image features and the user waits for that transfer
\cite{kang2017neurosurgeon}. Speculative decoding splits the work by
tokens: a small model drafts, a large model verifies, and the server still
reads the picture \cite{leviathan2023fast,chen2023accelerating}. Later
edge variants keep the same split
\cite{jia2026covspec,venkatesha2025fast}. edgeVLM answers on the device at
once and reuses a late cloud caption as context for later frames, rather
than as the result of the frame that is on screen now
\cite{qian2025edgevlm}. EdgeBoost uploads an input only when the local
confidence is low. On image classification this removes about 20--55\% of
the communication of a cloud-only run, and the cloud still returns a class
label \cite{said2024edgeboost}.

What these systems return is a class, a feature tensor, or a caption that
arrives after the frame has moved on. The prototype shows an activity
phrase as soon as the phone finishes, and it sends the picture only when
that phrase is unsafe relative to the action already on screen.

\subsection{Communication after the phrase exists}

A proportional-fair scheduler and an L4S queue assign service to packets
that have already been formed. They do not look at how long the laptop will
spend on the picture after it arrives
\cite{jalali2000data,briscoe2023l4s}. Semantic communication scores a
transmission by task success rather than by throughput, but the payload is
still a feature tensor \cite{xie2021deep}. Offloading chooses local or
remote execution, and adaptive streaming changes how many bytes are sent.
After the sentence has already been written, these mechanisms still leave
open whether the picture should follow, how large that picture should be,
and which phone should be checked first.

The proposal takes that last decision. Token confidence decides whether the
JPEG is uploaded. JPEG quality and resolution decide its size. For every
upload, the predicted edge latency sets the uplink deadline, and an
application-layer scheduler serves the earliest deadline first. The running
prototype implements the confidence decision and a fixed JPEG size. The
predictor and the scheduler are specified in Section~\ref{sec:design} and
are not in the code yet.
